% GENERATED FILE. Edit canonical lessons, then run npm run generate:books.
% canonical-source-hash: fnv1a64:baea2a0287b91858
% canonical-lessons: JA-C239-odoru, JA-C239-tsuru, JA-C239-sanpo, JA-C239-tejina, JA-C239-shougi

\chapter{Weekend Pastimes}
\label{ch:ja-weekend-pastimes}
\hlchaptermodality{239}

\begin{chapteropening}
\textbf{By the end of this chapter:} \emph{I can say how I spent a free weekend.}

Complete the last lesson of chapter 239: I can say how I spent a free weekend.
\end{chapteropening}

\section[\texorpdfstring{\ja{おどる} (odoru)}{おどる (odoru)}]{\ja{おどる} (odoru) — to dance}

\label{lesson:JA-C239-odoru}



\begin{quote}
Before the new one: say the Japanese for a programme, then the Japanese for boring.
\end{quote}

\subsection*{You'll want to know: \ja{おどる}}
\textbf{\ja{おどる}} — \emph{odoru} — \textquotedblleft{}to dance\textquotedblright{}.

\begin{quote}
\textbf{\ja{まつりで} \ja{みんなと} \ja{おどりました。}} — \emph{matsuri de minna to odorimashita} — \textquotedblleft{}I danced with everyone at the festival.\textquotedblright{}
\end{quote}

Polite \textbf{\ja{おどります}}, \emph{odorimasu}; a godan verb.

\begin{cousinweb}
\textbf{\ja{おどる}} (\emph{odoru}) is a native Japanese word (wago), inherited rather than borrowed from Chinese; its kanji was attached later as a kun'yomi (JMdict, KANJIDIC2).
\end{cousinweb}

\begin{grammarlens}[title={What you've built}]
One more word for everyday life at A2.
\end{grammarlens}

\subsection*{Your turn}
\begin{itemize}
\raggedright
  \item \emph{Say it:} \emph{odoru}
  \item \emph{Say it:} \emph{odoru}, once more
  \item \emph{Say it:} \emph{matsuri de minna to odorimashita} — I danced with everyone at the festival.
  \item \emph{Recall:} say the Japanese for a programme, then the Japanese for boring, then say \emph{odoru} again
\end{itemize}

\begin{tcolorbox}[breakable,colback=teal!4,colframe=teal!35!black,title={Before you move on}]
What does \ja{おどる} mean? (\textquotedblleft{}To dance\textquotedblright{}.) Say it once more.
\end{tcolorbox}
\section[\texorpdfstring{\ja{つる} (tsuru)}{つる (tsuru)}]{\ja{つる} (tsuru) — to go fishing; to catch (a fish)}

\label{lesson:JA-C239-tsuru}



\begin{quote}
Before the new one: say the Japanese for boring, then the Japanese for to dance.
\end{quote}

\subsection*{You'll want to know: \ja{つる}}
\textbf{\ja{つる}} — \emph{tsuru} — \textquotedblleft{}to go fishing; to catch (a fish)\textquotedblright{}.

\begin{quote}
\textbf{\ja{せんしゅう} \ja{かわで} \ja{さかなを} \ja{つりました。}} — \emph{senshū kawa de sakana o tsurimashita} — \textquotedblleft{}Last week I caught fish in the river.\textquotedblright{}
\end{quote}

Polite \textbf{\ja{つります}}, \emph{tsurimasu}; a godan verb.

\begin{cousinweb}
\textbf{\ja{つる}} (\emph{tsuru}) is a native Japanese word (wago), inherited rather than borrowed from Chinese; its kanji was attached later as a kun'yomi (JMdict, KANJIDIC2).
\end{cousinweb}

\begin{grammarlens}[title={What you've built}]
One more word for everyday life at A2.
\end{grammarlens}

\subsection*{Your turn}
\begin{itemize}
\raggedright
  \item \emph{Say it:} \emph{tsuru}
  \item \emph{Say it:} \emph{tsuru}, once more
  \item \emph{Say it:} \emph{senshū kawa de sakana o tsurimashita} — Last week I caught fish in the river.
  \item \emph{Recall:} say the Japanese for boring, then the Japanese for to dance, then say \emph{tsuru} again
\end{itemize}

\begin{tcolorbox}[breakable,colback=teal!4,colframe=teal!35!black,title={Before you move on}]
What does \ja{つる} mean? (\textquotedblleft{}To go fishing; to catch (a fish)\textquotedblright{}.) Say it once more.
\end{tcolorbox}
\section[\texorpdfstring{\ja{さんぽ} (sanpo)}{さんぽ (sanpo)}]{\ja{さんぽ} (sanpo) — a walk; a stroll}

\label{lesson:JA-C239-sanpo}



\begin{quote}
Before the new one: say the Japanese for to dance, then the Japanese for to go fishing.
\end{quote}

\subsection*{You'll want to know: \ja{さんぽ}}
\textbf{\ja{さんぽ}} — \emph{sanpo} — \textquotedblleft{}a walk; a stroll\textquotedblright{}.

\begin{quote}
\textbf{\ja{けさ} \ja{いぬと} \ja{こうえんを} \ja{さんぽしました。}} — \emph{kesa inu to kōen o sanpo shimashita} — \textquotedblleft{}This morning I took the dog for a walk in the park.\textquotedblright{}
\end{quote}

\begin{cousinweb}
\textbf{\ja{さんぽ}} (\emph{sanpo}) is Sino-Japanese (kango): two Chinese characters read with their on'yomi, readings borrowed from Chinese. They are \emph{san}, \textquotedblleft{}scatter\textquotedblright{} (Middle Chinese \emph{sanX}) and \emph{po}, \textquotedblleft{}step, walk\textquotedblright{} (\emph{buH}) (KANJIDIC2; Middle Chinese in Baxter's transcription).
\end{cousinweb}

\begin{grammarlens}[title={What you've built}]
One more word for everyday life at A2.
\end{grammarlens}

\subsection*{Your turn}
\begin{itemize}
\raggedright
  \item \emph{Say it:} \emph{sanpo}
  \item \emph{Say it:} \emph{sanpo}, once more
  \item \emph{Say it:} \emph{kesa inu to kōen o sanpo shimashita} — This morning I took the dog for a walk in the park.
  \item \emph{Recall:} say the Japanese for to dance, then the Japanese for to go fishing, then say \emph{sanpo} again
\end{itemize}

\begin{tcolorbox}[breakable,colback=teal!4,colframe=teal!35!black,title={Before you move on}]
What does \ja{さんぽ} mean? (\textquotedblleft{}A walk; a stroll\textquotedblright{}.) Say it once more.
\end{tcolorbox}
\section[\texorpdfstring{\ja{てじな} (tejina)}{てじな (tejina)}]{\ja{てじな} (tejina) — a magic trick; conjuring}

\label{lesson:JA-C239-tejina}



\begin{quote}
Before the new one: say the Japanese for to go fishing, then the Japanese for a walk.
\end{quote}

\subsection*{You'll want to know: \ja{てじな}}
\textbf{\ja{てじな}} — \emph{tejina} — \textquotedblleft{}a magic trick; conjuring\textquotedblright{}.

\begin{quote}
\textbf{\ja{おじさんは} \ja{まつりで} \ja{てじなを} \ja{しました。}} — \emph{ojisan wa matsuri de tejina o shimashita} — \textquotedblleft{}My uncle did magic tricks at the festival.\textquotedblright{}
\end{quote}

\begin{cousinweb}
\textbf{\ja{てじな}} (\emph{tejina}) is a native Japanese compound (wago), inherited rather than borrowed: it is built from \emph{te}, \textquotedblleft{}hand\textquotedblright{}, and \emph{shina}, \textquotedblleft{}goods, article\textquotedblright{}, with \emph{sh} voiced to \emph{j} (JMdict, KANJIDIC2).
\end{cousinweb}

\begin{grammarlens}[title={What you've built}]
One more word for everyday life at A2.
\end{grammarlens}

\subsection*{Your turn}
\begin{itemize}
\raggedright
  \item \emph{Say it:} \emph{tejina}
  \item \emph{Say it:} \emph{tejina}, once more
  \item \emph{Say it:} \emph{ojisan wa matsuri de tejina o shimashita} — My uncle did magic tricks at the festival.
  \item \emph{Recall:} say the Japanese for to go fishing, then the Japanese for a walk, then say \emph{tejina} again
\end{itemize}

\begin{tcolorbox}[breakable,colback=teal!4,colframe=teal!35!black,title={Before you move on}]
What does \ja{てじな} mean? (\textquotedblleft{}A magic trick; conjuring\textquotedblright{}.) Say it once more.
\end{tcolorbox}
\section[\texorpdfstring{\ja{しょうぎ} (shōgi)}{しょうぎ (shōgi)}]{\ja{しょうぎ} (shōgi) — shogi; Japanese chess}

\label{lesson:JA-C239-shougi}



\begin{quote}
Before the new one: say the Japanese for a walk, then the Japanese for a magic trick.
\end{quote}

\subsection*{You'll want to know: \ja{しょうぎ}}
\textbf{\ja{しょうぎ}} — \emph{shōgi} — \textquotedblleft{}shogi; Japanese chess\textquotedblright{}.

\begin{quote}
\textbf{\ja{きのう} \ja{ちちと} \ja{しょうぎを} \ja{しました。}} — \emph{kinō chichi to shōgi o shimashita} — \textquotedblleft{}Yesterday I played shogi with my father.\textquotedblright{}
\end{quote}

\begin{cousinweb}
\textbf{\ja{しょうぎ}} (\emph{shōgi}) is Sino-Japanese (kango): two Chinese characters read with their on'yomi, readings borrowed from Chinese. They are \emph{shō}, \textquotedblleft{}general\textquotedblright{} (Middle Chinese \emph{tsjangH}) and \emph{gi}, \textquotedblleft{}board game\textquotedblright{} (\emph{gi}) (KANJIDIC2; Middle Chinese in Baxter's transcription).
\end{cousinweb}

\begin{grammarlens}[title={What you've built}]
One more word for everyday life at A2.
\end{grammarlens}

\subsection*{Your turn}
\begin{itemize}
\raggedright
  \item \emph{Say it:} \emph{shōgi}
  \item \emph{Say it:} \emph{shōgi}, once more
  \item \emph{Say it:} \emph{kinō chichi to shōgi o shimashita} — Yesterday I played shogi with my father.
  \item \emph{Recall:} say the Japanese for a walk, then the Japanese for a magic trick, then say \emph{shōgi} again
\end{itemize}

\begin{tcolorbox}[breakable,colback=teal!4,colframe=teal!35!black,title={Before you move on}]
What does \ja{しょうぎ} mean? (\textquotedblleft{}Shogi; Japanese chess\textquotedblright{}.) Say it once more.
\end{tcolorbox}
